% =============================================================================
%  CHAPTER 3 - MATERIALS AND METHODOLOGY
% =============================================================================
\chapter[Methodology]{Materials and Methodology}
\label{chap:methodology}

\chapterquote{If you cannot weigh, measure, or photograph it, you cannot
yet engineer it.}{after Lord Kelvin, 1883}

\begin{chapterpreview}
This chapter documents in reproducible detail the apparatus, the
materials, and the protocol used in the experimental programme.
\cref{sec:meth-materials} characterises the basalt ballast and verifies
its conformity to the \gls{RDSO} gradation envelope.
\cref{sec:meth-PUBM,sec:meth-RUBM} specify the two under-ballast mat
materials. \cref{sec:meth-apparatus} describes the modified large-scale
drop-weight impact apparatus (Figure~\ref{fig:meth-flowchart}).
\cref{sec:meth-calibration} calibrates the apparatus against the
Indian Railways field-stress envelope. \cref{sec:meth-protocol} fixes
the assembly and blow protocol; \cref{sec:meth-matrix} sets out the
nine-cell factorial test matrix (Figure~\ref{fig:meth-matrix}). The
chapter ends with the twenty-two-descriptor per-particle pipeline
(\cref{sec:meth-morphology,sec:meth-descriptors}).
\end{chapterpreview}

\dropcap{T}{his} chapter documents the materials, the apparatus, the test protocol, and
the per-particle morphological-analysis pipeline used in the experimental
programme. The chapter is the longest in the thesis because the programme
combines two methodological strands---bulk-index testing through a modified
drop-weight impact apparatus, and paired three-dimensional morphological
characterisation of individually marked tracer particles---and
reproducibility of either strand requires a fully specified procedure. The
chapter is organised into eight sections following the natural workflow of
the programme: materials (Section~\ref{sec:meth-materials}), pad and base
conditions (Section~\ref{sec:meth-reinforcement}), tracer-particle
preparation (Section~\ref{sec:meth-scanning}), the modified drop-weight
impact apparatus (Section~\ref{sec:meth-apparatus}), stress-and-energy
calibration to Indian Railways field loading
(Section~\ref{sec:meth-calibration}), the blow-by-blow test protocol and
bulk-index computation (Section~\ref{sec:meth-protocol}), and the image-based
three-dimensional morphological-analysis pipeline
(Section~\ref{sec:meth-morphology}).

\begin{figure}[ht]
\centering
\begin{tikzpicture}[node distance=0.5cm and 0.6cm,>=Stealth]
    % -------- Bulk stream (left column) --------
    \node[flowstart]                                            (start)    {Source basalt ballast (RDSO gradation)};
    \node[flowstep, below=of start]                             (sieve-in) {Dry-sieve analysis\\(pre-test PSD)};
    \node[flowstep, below=of sieve-in]                          (prep)     {Specimen assembly\\(mould + base + pad + 20\,kg fill)};
    \node[flowstep, below=of prep]                              (impact)   {Drop-weight impact test\\(20 blows; 300 or 450\,mm drop)};
    \node[flowstep, below=of impact]                            (sett)     {Settlement \& peak force\\(every 5 blows)};
    \node[flowstep, below=of sett]                              (sieve-out){Post-test sieve analysis\\(RDSO series)};
    \node[flowstep, below=of sieve-out]                         (fiin)     {Compute $\FIIN = \PoneM + \PtwoM$};

    % -------- Per-particle stream (right column) --------
    \node[flowstep, right=2.6cm of prep]                        (select)   {Select 5 tracer particles\\per configuration ($n=45$)};
    \node[flowstep, below=of select]                            (scanbef)  {Pre-test 3D scan\\(EinScan SL, $\sim$0.1\,mm)};
    \node[flowstep, below=of scanbef]                           (qcbef)    {QC: watertight,\\single body, vol.\ in band};
    \node[flowstep, below=of qcbef]                             (extract)  {Recover tracers post-test,\\re-scan, QC};
    \node[flowstep, below=of extract]                           (desc)     {Compute 22 descriptors\\and paired $\Delta X$, $\%\Delta X$};

    % Synthesis box
    \node[flowend, below=2.2cm of $(fiin)!0.5!(desc)$]          (synth)    {Synthesis (\FIIN{} + paired morphology)};

    % -------- Arrows --------
    \draw[flowarrow] (start)    -- (sieve-in);
    \draw[flowarrow] (sieve-in) -- (prep);
    \draw[flowarrow] (prep)     -- (impact);
    \draw[flowarrow] (impact)   -- (sett);
    \draw[flowarrow] (sett)     -- (sieve-out);
    \draw[flowarrow] (sieve-out)-- (fiin);
    \draw[flowarrow] (fiin)     |- (synth);

    \draw[flowarrow] (select)   -- (scanbef);
    \draw[flowarrow] (scanbef)  -- (qcbef);
    \draw[flowarrow, dashed]    (qcbef.south)
        .. controls +(0,-0.4) and +(2.5,0) ..
        (prep.east) node[midway,above,font=\scriptsize\itshape] {tracers go in};
    \draw[flowarrow, dashed]    (impact.east)
        .. controls +(2.5,0) and +(0,0.4) ..
        (extract.north) node[midway,below,font=\scriptsize\itshape] {tracers come out};
    \draw[flowarrow] (extract)  -- (desc);
    \draw[flowarrow] (desc)     |- (synth);

    % Stream labels
    \node[font=\small\bfseries\color{gsvblue}, above=0.2cm of start]
        {Bulk stream};
    \node[font=\small\bfseries\color{gsvblue}, above=0.2cm of select]
        {Per-particle stream};
\end{tikzpicture}
\caption{Experimental methodology: a two-stream pipeline. The bulk stream
         (left) follows the specimen through assembly, impact testing,
         and post-test sieve analysis, ending in the Indian Railway
         Fouling Index \FIIN. The per-particle stream (right) follows 45
         individually-tracked tracer particles through 3D scanning,
         quality control, impact loading, re-scan, and 22-descriptor
         comparison. The two streams converge in the synthesis stage.}
\label{fig:meth-flowchart}
\end{figure}

\section{Basalt ballast}
\label{sec:meth-materials}

\subsection{Source and parent rock}
\label{sec:meth-source}

All ballast used in the experimental programme was sourced from a single
quarry supplying \gls{RDSO}-specification railway ballast in the Vadodara
region of Gujarat, India. A single source was used to eliminate parent-rock
variability as a confounding factor between tests: any cross-test difference
in degradation can then be attributed to the reinforcement configuration
rather than to unrecorded mineralogical variation. The parent rock is
basalt---a fine-grained, dense, aphanitic igneous rock with a dry particle
density \SI{2.76}{\gram\per\centi\metre\cubed} and low water absorption
(\(<\SI{1}{\percent}\)). These figures are consistent with the \gls{RDSO}
specification \citep{RDSO2023} requirement of not less than
\SI{2.60}{\gram\per\centi\metre\cubed} dry density and with the parent-rock
properties reported for comparable basalt in the impact literature
\citep{Koohmishi2020, Zhang2023}.

Approximately \SI{25000}{\kilogram} of ballast was procured for the
programme---much more than required for the nine main tests (approximately
\SI{20}{\kilogram} each), for the preliminary calibration trials run during
apparatus commissioning, and for a replacement reserve. All nine tests
therefore drew from the same stockpile and the same parent rock. The
stockpile was kept covered and elevated from direct ground contact between
tests to prevent contamination of the void-fouling baseline.

\subsection{Particle size distribution and the RDSO envelope}
\label{sec:meth-psd}

The as-supplied ballast was characterised by dry-sieve analysis using the
following sieve series: 63, 50, 40, 31.5, 25, 20, 16, 12.5, 10, 4.75, 2.36,
0.6, 0.15 and \SI{0.075}{\milli\metre}. Sieving was performed on sample
masses of approximately \SI{10}{\kilogram} in accordance with the spirit of
the \gls{RDSO} procedure \citep{RDSO2023}. The resulting gradation, plotted
in \cref{fig:meth-psd} against the \gls{RDSO} upper and lower envelope
curves, shows that the ballast conforms to the specification with a small
reserve against both limits---the target gradation sits close to the mid-band
of the envelope throughout the \SIrange{20}{63}{\milli\metre} range. The
percentage retained on the \SI{65}{\milli\metre} sieve was below the
\SI{5}{\percent} maximum permitted by the specification; the percentage
passing the \SI{20}{\milli\metre} sieve was below \SI{2}{\percent}, well
within the \SI{2}{\percent} maximum fines limit. The gradation is therefore
centred within the acceptable window and is representative of
production-quality \gls{RDSO} ballast.

\begin{figure}[ht]
    \centering
    \includegraphics[width=0.85\textwidth]{image10.png}
    \caption{Pre-test particle size distribution of the basalt ballast used
             in the experimental programme, plotted against the upper and
             lower envelope curves of the \acrshort{RDSO} specification.}
    \label{fig:meth-psd}
\end{figure}

\subsection{Physical and mechanical properties}
\label{sec:meth-properties}

The ballast was characterised for the physical and mechanical properties
required by the \gls{RDSO} specification. The measurements are summarised in
\cref{tab:meth-ballast-properties}. All values fall within the \gls{RDSO}
acceptance thresholds; the \gls{AIV} measured on the artificially crushed
\SIrange{10}{12.5}{\milli\metre} sub-sample is reported here for
completeness but is precisely the index whose limitations are set out in
Chapters~\ref{chap:introduction} and~\ref{chap:literature}. The remainder of
the programme accordingly reports the bulk fouling index \FIIN{} on the
full-size specimen.

\begin{table}[ht]
\centering
\caption{Physical and mechanical properties of the basalt ballast used in
         the experimental programme.}
\label{tab:meth-ballast-properties}
\rowcolors{2}{rowbg}{white}
\begin{tabular}{@{}lccl@{}}
    \toprule
    \rowcolor{tablehdr}
    {\color{white}\textbf{Property}}
    & {\color{white}\textbf{Measured}}
    & {\color{white}\textbf{\acrshort{RDSO} limit}}
    & {\color{white}\textbf{Method}}\\
    \midrule
    Dry particle density (\si{\gram\per\centi\metre\cubed})
        & 2.78 & \(\ge 2.60\) & Gas pycnometer $+$ weighing\\
    Bulk (loose) density (\si{\kilogram\per\metre\cubed})
        & 1460 & Not specified & Bulk volume in mould\\
    Water absorption (\si{\percent})
        & 0.72 & \(\le 1.0\) & ASTM~C127 equivalent\\
    \acrshort{AAV} (\acrshort{LAA}) (\si{\percent})
        & 15.4 & \(\le 30\) & IS:2386 Pt~IV\\
    \acrshort{AIV} (\si{\percent})
        & 12.1 & \(\le 20\) & IS:2386 Pt~IV (conventional)\\
    \bottomrule
\end{tabular}
\end{table}

\begin{figure}[ht]
    \centering
    \includegraphics[width=0.70\textwidth]{image11.jpeg}
    \caption{Pile of the \acrshort{RDSO}-gradation basalt ballast used in
             the experimental programme (Civil Lab, \acrshort{GSV}).}
    \label{fig:meth-ballast-pile}
\end{figure}

\section{Reinforcement Conditions}
\label{sec:meth-reinforcement}

Three pad conditions and two base conditions are crossed in the factorial
test matrix. The three pad conditions are an Unreinforced (\gls{UR})
control, a Rubber Under-Ballast Mat (\gls{RUBM}), and a Polyurethane
Under-Ballast Mat (\gls{PUBM}). The two base conditions are a Flexible Base
(\gls{FB}) of compacted sand and a Rigid Base (\gls{RB}) of the concrete
mould floor. The systematic three-letter abbreviations are introduced
formally in this thesis as a replacement for the ad-hoc numbered test
labels used in earlier laboratory work.

\subsection{Unreinforced (UR) control}
\label{sec:meth-UR}

The \gls{UR} configuration places no pad between the ballast fill and the
base element of the mould. The \gls{UR} condition serves three analytical
purposes simultaneously: it is the mechanical control against which the two
pad conditions are compared; it represents the in-service condition of
legacy ballasted track with no substructure reinforcement; and it isolates
the effect of the base condition (flexible sand versus rigid concrete) from
the effect of any compliant pad placed above it. The \gls{UR} condition is
present in three of the nine test configurations, providing the baseline
against which the pad-protected configurations are measured.

\subsection{Rubber Under-Ballast Mat (RUBM)}
\label{sec:meth-RUBM}

The \gls{RUBM} material is a commercially manufactured \SI{12}{\milli\metre}
thick sheet of vulcanised natural rubber with a Shore~A hardness of
approximately~65 and a nominal bulk density of
\SI{1.35}{\gram\per\centi\metre\cubed}. The sheet was cut to a circular disc
matching the \SI{300}{\milli\metre} internal diameter of the test mould and
placed flat on the base element before the ballast fill. The compliance of
the \gls{RUBM} falls in the middle of the range reported in the literature
for natural-rubber under-ballast mats
\citep{Navaratnarajah2017, Ngo2019}---sufficient to dissipate a meaningful
fraction of the hammer impulse energy while stiff enough to recover
elastically between blows without progressive plastic deformation of the pad
itself. No \gls{RUBM} specimen was reused between tests; a fresh
\SI{12}{\milli\metre} disc was cut for every test to eliminate any
carry-over effect of cumulative pad damage.

\begin{figure}[ht]
    \centering
    \includegraphics[width=0.55\textwidth]{image12.jpeg}
    \caption{Rubber Under-Ballast Mat (\acrshort{RUBM}) disc,
             \SI{12}{\milli\metre} thick natural rubber, cut to
             \SI{300}{\milli\metre} diameter.}
    \label{fig:meth-RUBM}
\end{figure}

\subsection{Polyurethane Under-Ballast Mat (PUBM)}
\label{sec:meth-PUBM}

The \gls{PUBM} material is a commercial Getzner polyurethane under-ballast
mat product. Getzner is a widely adopted supplier of polyurethane
track-resilient elements on European, Australian, and Indian mainline and
urban rail networks, and the specific variant used in this study is
representative of the stiffness range in which Getzner UBM products are
typically deployed beneath ballast in transition zones and impact-prone
locations. The \gls{PUBM} disc was cut to \SI{300}{\milli\metre} diameter
matching the mould internal diameter, like the \gls{RUBM}. \citet{Gundavaram2023}
have previously demonstrated, in cyclic-triaxial tests on Indian ballast,
that an elastomeric polyurethane inclusion materially reduces both
cumulative plastic deformation and the rate of fines generation, and the
use of \gls{PUBM} in the present programme extends that line of evidence
from cyclic to impact loading.

\begin{figure}[ht]
    \centering
    \includegraphics[width=0.55\textwidth]{image13.jpeg}
    \caption{Polyurethane Under-Ballast Mat (\acrshort{PUBM}, commercial
             Getzner product), cut to \SI{300}{\milli\metre} diameter.}
    \label{fig:meth-PUBM}
\end{figure}

\subsection{Base conditions (FB / RB)}
\label{sec:meth-base}

The base condition---the element directly beneath the pad, or beneath the
ballast if no pad is used---is the second factorial variable in the
experimental design. The \gls{FB} condition comprises a
\SI{100}{\milli\metre} deep cushion of clean, medium-grained quartz sand
placed at the base of the test mould before the pad and ballast fill. The
sand was oven-dried to constant mass and compacted in the mould by a
standardised vibratory compaction protocol to a bulk density of
approximately \SI{1.60}{\gram\per\centi\metre\cubed}. The \gls{RB}
condition places the ballast fill (with or without the interposed pad)
directly on the concrete floor of the test mould. The concrete floor is
several orders of magnitude stiffer than either the pad or the sand base
and is effectively unyielding at the stress levels delivered by the hammer.

The two base conditions are chosen to bracket the stiffness range that a
ballast layer might encounter in service: the \gls{FB} approximates a
well-compacted granular sub-ballast, while the \gls{RB} approximates an
unyielding structural element such as a bridge deck, a concrete track-slab
foundation, or a poorly-maintained subgrade that has densified below its
own elastic limit. The factorial crossing of pad and base conditions allows
the pad-effect contribution and the base-effect contribution to be
separated cleanly, which is the analytical design that \citet{Nimbalkar2012}
identified as essential for distinguishing pad performance from sub-grade
compliance in laboratory shock-mat work.

\begin{figure}[ht]
    \centering
    \includegraphics[width=0.55\textwidth]{image14.jpeg}
    \caption{Test mould prepared with the \SI{100}{\milli\metre} flexible
             base sand cushion (\acrshort{FB} condition), photographed
             immediately before placement of the pad and the ballast fill.}
    \label{fig:meth-FB}
\end{figure}

\section{Particle scanning}
\label{sec:meth-scanning}

\subsection{Rationale for the paired-particle design}
\label{sec:meth-paired-rationale}

The thesis rests on a paired-particle experimental design: the same physical
grains are scanned before and after the impact-loading event, and the
grain-scale damage is tracked within each particle rather than averaged
across the population.

\begin{mechanism}[title=Why pair before and after?]
Pairing eliminates the particle-to-particle variability in initial shape
---which in any typical ballast assembly is considerable, with volumes
ranging over almost an order of magnitude across the
\SIrange{20}{65}{\milli\metre} gradation---as a source of scatter in the
damage signal. It also enables the direct attribution of each observed
change in shape descriptors to the loading event, because any change
recorded in the after-scan that is not also present in the before-scan of
the same particle must have arisen from the hammer blows.
\end{mechanism}

\subsection{Selection criteria}
\label{sec:meth-selection}

Tracer particles were selected from the bulk stockpile by visual inspection
and by calliper measurement. Four criteria were applied in sequence; a
candidate particle had to satisfy all four to be accepted as a tracer.

\begin{enumerate}[label=(\roman*),leftmargin=*]
    \item \textbf{Size:} the equivalent spherical diameter---estimated at
          the selection stage from the three calliper-measured principal
          dimensions---had to fall in the
          \SIrange{30}{65}{\milli\metre} range.
    \item \textbf{Shape:} particles with extreme flatness or elongation
          (aspect ratio $L/S>3$ by initial calliper reading) were rejected
          because such shapes produce preferentially high breakage under
          impact loading and would distort the \gls{ModeA} / \gls{ModeB}
          distribution.
    \item \textbf{Parent freshness:} candidates showing any visible
          weathering rind, iron-oxide staining, or surface cementation were
          rejected.
    \item \textbf{Identifiability:} the selected particles had to admit a
          unique, durable surface marking for unambiguous before--after
          pairing.
\end{enumerate}

Each selected tracer particle was individually marked with a unique code.
The marking protocol uses a paint marker applied to a clean, dry surface
patch on each particle. The paint coat is thin
(\(<\SI{0.2}{\milli\metre}\)) and does not measurably affect either the
volume or the surface area computed from the scans.

\begin{figure}[ht]
    \centering
    \includegraphics[width=0.65\textwidth]{image15.jpeg}
    \caption{Marked tracer particles showing the white acrylic paint patch
             and the hand-written numeric identifier code.}
    \label{fig:meth-tracers}
\end{figure}

\subsection{Pre-test structured-light scanning}
\label{sec:meth-prescan}

Each tracer particle was scanned on an EinScan (Shining 3D Technology
Co.\ Ltd., Hangzhou, China) desktop structured-light scanner. The
EinScan~SE was used across the programme, providing comparable geometric
accuracy of approximately \SI{0.1}{\milli\metre} and scan resolution of
approximately \SI{0.15}{\milli\metre} in the automatic turntable mode used
here. The detailed scan-protocol and mesh quality-control procedures are
set out in \cref{sec:meth-morphology} below, in the context of the
morphological-analysis pipeline. In the present sub-section pre-test
scanning is described as part of specimen preparation: every tracer
particle carried a complete, validated, watertight triangular surface mesh
in \gls{STL} format into the impact test, and the pre-test mesh serves as
the reference state against which the post-test mesh is paired.

\begin{figure}[ht]
    \centering
    \includegraphics[width=0.65\textwidth]{image16.jpeg}
    \caption{EinScan structured-light scanner (Shining~3D) in operation with
             a marked tracer particle mounted on the automatic turntable.}
    \label{fig:meth-einscan}
\end{figure}

\section{The modified drop-weight impact apparatus}
\label{sec:meth-apparatus}

The apparatus comprises five principal components: a cylindrical cast-iron
test mould of \SI{300}{\milli\metre} internal diameter and
\SI{350}{\milli\metre} internal depth, seated on a rigid reaction
base-plate; a vertical guide frame of four mild-steel columns connected by
an upper crosshead; a cylindrical cast-iron hammer of \SI{50}{\kilogram}
mass, sliding on the guide columns in free fall; a continuously adjustable
hammer release-and-locking mechanism graduated over the
\SIrange{200}{450}{\milli\metre} range of drop height; and a load-transfer
assembly consisting of a \SI{300}{\milli\metre} bearing plate, a
compression-type \SI{5}{\tonne} strain-gauge load cell, and a hardened-steel
striker plate. A mechanical blow counter and two manual operating handles
complete the assembly. The mould diameter of \SI{300}{\milli\metre} delivers
a $D/\dmax$ ratio of approximately~6 against the \gls{RDSO}
\SI{50}{\milli\metre} mean particle size, eliminating the boundary-dominance
of the conventional IS:2386 Part~IV apparatus \citep{BIS1963} and bringing
the present rig into the design envelope established by
\citet{Nimbalkar2012} and the subsequent impact literature surveyed in
\cref{sec:lit-lab}.

The adjustable release mechanism permits the hammer to be locked at any
drop height between \SI{200}{\milli\metre} and \SI{450}{\milli\metre}
against a graduated scale on one of the four guide columns, so that the
same apparatus can deliver a range of impact energies without disassembly.
The in-line load cell measures the dynamic peak force imparted through the
striker plate on every blow and transmits the signal to a digital
peak-force indicator; this allows each test to be characterised not only by
its nominal drop-energy but by the actual force time-series, which is the
principal advance of the present rig over the conventional apparatus. The
mechanical blow counter provides a redundant and power-independent count of
delivered blows. Detailed design drawings, material specifications,
tolerance callouts, and the safety-interlock documentation are set out in
the utility patent application filed for this apparatus.

\begin{figure}[ht]
\centering
\begin{tikzpicture}[scale=0.9,>=Stealth,every node/.style={font=\footnotesize}]
    % --- Frame columns (4 guide rods) ---
    \draw[gsvblue!40,line width=1.5pt] (-3.2,-3.2) -- (-3.2,7.2);
    \draw[gsvblue!40,line width=1.5pt] ( 3.2,-3.2) -- ( 3.2,7.2);
    \node[font=\scriptsize,gsvblue,rotate=90] at (-3.7,2)
        {\bfseries Guide column (4)};

    % --- Top crosshead ---
    \fill[gsvblue!80] (-3.4, 7.0) rectangle (3.4, 7.4);
    \draw[gsvblue,thick] (-3.4, 7.0) rectangle (3.4, 7.4);
    \node[white,font=\scriptsize\bfseries] at (0,7.2) {Top crosshead};

    % --- Release / lifting mechanism ---
    \draw[gsvaccent,thick,fill=gsvaccent!20] (-0.5,6.3) rectangle (0.5,7.0);
    \node[font=\tiny,gsvaccent!80!black] at (0,6.65) {release};
    \node[font=\scriptsize,gsvblue!90!black,anchor=west]
        at (0.8,6.65) {Mechanical release latch};

    % --- Hammer (50 kg) ---
    \draw[gsvblue,thick,fill=gsvblue!85,rounded corners=2pt]
        (-1.7,4.6) rectangle (1.7,5.7);
    \node[white,font=\bfseries\small] at (0,5.15) {Hammer (50\,kg)};
    % Drop arrow
    \draw[->,line width=1.4pt,red!70] (0,4.55) -- (0,3.65)
        node[midway,right=4pt,red!80,font=\scriptsize] {drop $h$};
    % Drop-height markings
    \draw[gsvblue!50,dotted] (1.7,5.15) -- (3.3,5.15);
    \node[font=\tiny,gsvblue!70,anchor=west] at (3.3,5.15)
        {hammer release elevation};

    % --- Striker plate ---
    \draw[gray!60,thick,fill=gray!40] (-1.5,3.3) rectangle (1.5,3.55);
    \node[font=\scriptsize,gsvblue!90!black,anchor=west] at (1.7,3.43)
        {Striker plate (\SI{300}{\milli\metre})};

    % --- In-line load cell ---
    \draw[gsvaccent!80!black,thick,fill=gsvaccent!30]
        (-0.85,3.0) rectangle (0.85,3.3);
    \node[font=\tiny,gsvaccent!80!black] at (0,3.15) {load cell};
    \node[font=\scriptsize,gsvblue!90!black,anchor=west] at (1.7,3.15)
        {In-line \SI{5}{\tonne} load cell};

    % --- Bearing plate ---
    \draw[gray!60,thick,fill=gray!40] (-1.6,2.7) rectangle (1.6,3.0);
    \node[font=\scriptsize,gsvblue!90!black,anchor=west] at (1.7,2.85)
        {Bearing plate};

    % --- Mould (cylindrical specimen container) ---
    \draw[gsvblue,thick,fill=gsvblue!8]
        (-2.3,0.5) -- (-2.3,2.7) -- (-2.0,2.7) -- (-2.0,0.5)
        -- (2.0,0.5) -- (2.0,2.7) -- (2.3,2.7) -- (2.3,0.5)
        -- cycle;
    % Specimen contents -- ballast layer
    \begin{scope}[clip]
        \path (-2.0,0.9) rectangle (2.0,2.65);
    \end{scope}
    \foreach \x/\y/\r in {
        -1.6/2.4/0.16, -1.1/2.45/0.18, -0.4/2.4/0.15, 0.3/2.5/0.17, 1.0/2.4/0.16, 1.55/2.45/0.16,
        -1.7/2.0/0.18, -1.0/2.05/0.16, -0.2/2.0/0.18, 0.5/2.1/0.16, 1.2/2.0/0.17, 1.7/2.05/0.15,
        -1.5/1.6/0.17, -0.7/1.55/0.18, 0.0/1.6/0.16, 0.7/1.55/0.18, 1.4/1.6/0.16,
        -1.7/1.15/0.16, -1.0/1.2/0.17, -0.2/1.15/0.16, 0.5/1.2/0.18, 1.2/1.15/0.16, 1.7/1.2/0.15}
        \fill[gray!55] (\x,\y) circle (\r);
    \node[font=\scriptsize,gsvblue!90!black,anchor=west] at (2.6,2.0)
        {Ballast fill (20\,kg)};
    % Pad layer
    \fill[orange!50] (-2.0,0.85) rectangle (2.0,0.95);
    \node[font=\scriptsize,orange!80!black,anchor=west] at (2.6,0.9)
        {Pad (PUBM/RUBM/UR)};
    % Base layer
    \fill[gsvblue!90] (-2.0,0.55) rectangle (2.0,0.85);
    \node[font=\scriptsize,gsvblue!90!black,anchor=west] at (2.6,0.7)
        {Base (FB / RB)};

    % --- Floor / base of apparatus ---
    \fill[gray!50] (-3.4,-0.2) rectangle (3.4,0.5);
    \draw[gsvblue,thick] (-3.4,-0.2) rectangle (3.4,0.5);
    \node[white,font=\scriptsize\bfseries] at (0,0.15) {Steel base plate};

    % --- Concrete floor (hatched) ---
    \fill[gray!30,pattern=north east lines] (-3.4,-1.5) rectangle (3.4,-0.2);
    \node[font=\scriptsize,gray!50!black,anchor=west] at (-3.0,-0.85)
        {Reinforced concrete floor};

    % --- Dimensional callouts ---
    \draw[<->,gsvblue!70] (-3.6, 5.15) -- (-3.6, 3.55)
        node[midway,left,gsvblue,font=\scriptsize,align=right]
        {200--450\,mm\\drop range};

    % --- Digital indicator ---
    \draw[gsvblue,thick,fill=white]
        (3.9, 4.0) rectangle (5.1, 5.0);
    \node[font=\scriptsize\bfseries,gsvblue] at (4.5, 4.7) {LCD};
    \node[font=\tiny,red!80,font=\ttfamily] at (4.5, 4.4) {37.0\,kN};
    \node[font=\tiny,gsvblue] at (4.5, 4.15) {peak force};
    \draw[->,gsvblue!60,thick] (3.6, 3.15) -- (3.9, 4.0);
\end{tikzpicture}
\caption{Schematic of the modified large-scale drop-weight impact
         apparatus. The hammer (\SI{50}{\kilogram}) is locked at a
         release elevation between \SI{200}{\milli\metre} and
         \SI{450}{\milli\metre} above the striker plate. The in-line
         \SI{5}{\tonne} load cell records the peak contact force at
         every blow and displays the value on the LCD indicator (right).
         The specimen mould holds the layered base, pad, and
         \SI{20}{\kilogram} ballast fill. Four guide columns ensure
         vertical impact alignment.}
\label{fig:meth-apparatus-schematic}
\end{figure}

\begin{figure}[ht]
    \centering
    \begin{subfigure}{0.46\textwidth}
        \centering
        \includegraphics[width=\linewidth]{image17.png}
        \caption{Full apparatus view.}
    \end{subfigure}\hfill
    \begin{subfigure}{0.46\textwidth}
        \centering
        \includegraphics[width=\linewidth]{image18.png}
        \caption{Load-transfer assembly.}
    \end{subfigure}
    \caption{Modified drop-weight impact apparatus (photographs
             matching Figure~\ref{fig:meth-apparatus-schematic}). (a)
             Overall view of the apparatus and its frame. (b) Close-up
             of the load-transfer assembly (bearing plate,
             \SI{5}{\tonne} load cell, striker plate) and the digital
             peak-force display.}
    \label{fig:meth-apparatus}
\end{figure}

\subsection{Comparison with the conventional AIV apparatus}
\label{sec:meth-AIV-comparison}

The salient design differences between the conventional IS:2386 Part~IV
\gls{AIV} apparatus and the modified large-scale apparatus commissioned at
\gls{GSV} are tabulated in \cref{tab:meth-AIV-comparison}.

\begin{table}[ht]
\centering
\caption{Principal specification of the modified apparatus in comparison to
         the conventional IS:2386 Part~IV Aggregate Impact Value apparatus.
         Only the headline parameters are listed; the full specification is
         in the patent documentation (\cref{chap:appB}).}
\label{tab:meth-AIV-comparison}
\small
\rowcolors{2}{rowbg}{white}
\begin{adjustbox}{max width=\textwidth}
\begin{tabular}{@{}p{4.4cm}p{4.8cm}p{5.6cm}@{}}
    \toprule
    \rowcolor{tablehdr}
    {\color{white}\textbf{Parameter}}
    & {\color{white}\textbf{Conventional \acrshort{AIV} (IS:2386 Pt~IV)}}
    & {\color{white}\textbf{Modified apparatus (present)}}\\
    \midrule
    Mould internal diameter
        & \SI{102}{\milli\metre} & \SI{300}{\milli\metre}\\
    Mould effective depth
        & \(\sim\!\SI{50}{\milli\metre}\) & \SI{350}{\milli\metre}\\
    Hammer mass
        & \SI{13.5}{\kilogram} & \SI{50}{\kilogram}\\
    Drop height
        & \SI{380}{\milli\metre} (fixed)
        & \SIrange{200}{450}{\milli\metre} (adjustable)\\
    Particle size tested
        & \SIrange{10}{12.5}{\milli\metre} (crushed sub-sample)
        & \SIrange{20}{63}{\milli\metre} (full \acrshort{RDSO} gradation)\\
    Approx.\ particles per test
        & \numrange{2}{4} & \numrange{40}{80}\\
    $D/\dmax$ ratio
        & \(\approx 1.6\) (boundary-dominated)
        & \(\approx 6\) (boundary-free)\\
    Stress simulation
        & Fixed / none
        & \SIrange{275}{523}{\kilo\pascal} (adjustable, maps to
        \SI{22.5}{\tonne} axle)\\
    Real-time force measurement
        & None
        & \SI{5}{\tonne} load cell $+$ digital peak-force indicator\\
    Blow counting
        & Manual & Integrated mechanical counter\\
    Post-test degradation metric
        & \acrshort{AIV} (\(\%\) fines through \SI{2.36}{\milli\metre})
        & \FIIN{} $+$ paired per-particle morphology\\
    \bottomrule
\end{tabular}
\end{adjustbox}
\end{table}

\section{Stress and energy calibration to Indian Railways field loading}
\label{sec:meth-calibration}

The calibration of the apparatus to in-service stress conditions proceeds
in two steps: first, the in-service stress range at the top of the ballast
layer under Indian Railways axle loading is estimated; second, the hammer
mass and drop-height range of the apparatus are chosen to span this
in-service range when divided by the mould cross-sectional area.

\subsection{Field stress estimation}
\label{sec:meth-field-stress}

The wheel load is distributed through the rail to the sleeper beneath it,
and thence to the ballast over the effective sleeper--ballast bearing area.
For a pre-stressed concrete mainline sleeper on Indian Railways the
effective bearing area is approximately
\SIrange{2000}{2200}{\centi\metre\squared}, of which the two rail-seat
regions each carry roughly half the axle load on each rail. Taking $\Ab
\approx \SI{1100}{\centi\metre\squared}$ as the effective bearing area
below a single rail seat, the static contact stress at the top of the
ballast under a \SI{22.5}{\tonne} axle is
\begin{equation}
    \sigma_{\text{static}} = \frac{W}{\Ab}
    \approx \frac{\SI{110.25}{\kilo\newton}}{\SI{0.11}{\metre\squared}}
    \approx \SI{1000}{\kilo\newton\per\metre\squared}
    \approx \SI{250}{\kilo\pascal}.
    \label{eq:meth-sigma-static}
\end{equation}
Applying a conservative dynamic amplification of~1.5 under ordinary
mainline traffic yields $\sigma \approx \SI{375}{\kilo\pascal}$; at
impact-prone locations where the amplification can reach $2\times$ or
higher, the contact stress rises into the
\SIrange{500}{600}{\kilo\pascal} range. The in-service stress window at the
top of the ballast layer under Indian Railways traffic is therefore
conservatively taken as approximately \SIrange{250}{550}{\kilo\pascal},
with the upper end reflecting heavy-haul and transition-zone conditions.

\subsection{Apparatus stress and energy}
\label{sec:meth-apparatus-stress}

For the modified apparatus the nominal contact stress $\sigma$ is the peak
hammer force $F$ divided by the bearing-plate area $\Ap$. The hammer
kinetic energy at impact is
\begin{equation}
    \Ke = m\,g\,h,
    \label{eq:meth-ke}
\end{equation}
where $m$ is the hammer mass (\SI{50}{\kilogram}), $g$ is the gravitational
acceleration (\SI{9.81}{\metre\per\second\squared}), and $h$ is the drop
height. Using the measured peak forces reported at commissioning---$F
\approx \SI{19.4}{\kilo\newton}$ at $h = \SI{300}{\milli\metre}$ and $F
\approx \SI{37.0}{\kilo\newton}$ at $h = \SI{450}{\milli\metre}$---and the
bearing-plate area $\Ap = \pi\,(\SI{0.150}{\metre})^{2} =
\SI{0.0707}{\metre\squared}$, the two apparatus operating points correspond
to
\begin{align}
    h &= \SI{300}{\milli\metre}: &
    \Ke &= 50 \times 9.81 \times 0.300 = \SI{147}{\joule}, &
    \sigma &= \frac{\SI{19.4}{\kilo\newton}}{\SI{0.0707}{\metre\squared}}
        \approx \SI{275}{\kilo\pascal};
        \label{eq:meth-h300}\\
    h &= \SI{450}{\milli\metre}: &
    \Ke &= 50 \times 9.81 \times 0.450 = \SI{221}{\joule}, &
    \sigma &= \frac{\SI{37.0}{\kilo\newton}}{\SI{0.0707}{\metre\squared}}
        \approx \SI{523}{\kilo\pascal}.
        \label{eq:meth-h450}
\end{align}

\begin{finding}[title=Calibration outcome]
The two operating points---\SI{275}{\kilo\pascal} at \SI{147}{\joule} and
\SI{523}{\kilo\pascal} at \SI{221}{\joule}---correspond respectively to the
lower end (ordinary mainline traffic) and the upper end (heavy-haul and
transition-zone conditions) of the in-service stress window. The two drop
heights selected for the test programme therefore bracket the Indian
Railways field-loading envelope cleanly.
\end{finding}

\begin{figure}[ht]
\centering
\begin{tikzpicture}[>=Stealth,every node/.style={font=\footnotesize}]
    % Vertical stress axis
    \draw[thick,->] (0,0) -- (0,7.2) node[above]
        {\scriptsize $\sigma$ (\si{\kilo\pascal})};
    \foreach \y/\v in {0.5/250,1.5/300,2.5/350,3.5/400,4.5/450,5.5/500,6.5/550}
        \draw (-0.1,\y) -- (0.1,\y) node[left,xshift=-0.15cm] {\v};

    % Field stress band
    \fill[gsvblue!15] (0.7, 0.5) rectangle (4.5, 6.5);
    \node[gsvblue,align=center] at (2.6, 3.5)
        {In-service stress\\window\\\textbf{250--550 \si{\kilo\pascal}}};
    \node[gsvblue,anchor=west,font=\scriptsize\itshape] at (0.7, 0.3)
        {static \SI{250}{\kilo\pascal}};
    \node[gsvblue,anchor=west,font=\scriptsize\itshape] at (0.7, 6.7)
        {impact-zone \SI{550}{\kilo\pascal}};

    % Apparatus operating points
    \draw[->,thick,boxmech!80] (5.6, 0.75) -- (4.5, 0.75);
    \node[fill=boxmech,text=white,rounded corners=2pt,
          font=\small\bfseries] at (6.2, 0.75)
        {\SI{300}{\milli\metre} drop};
    \node[anchor=west,font=\footnotesize] at (7.2, 0.75)
        {$\sigma\approx\SI{275}{\kilo\pascal}$,\ $\Ke=\SI{147}{\joule}$};

    \draw[->,thick,boxmech!80] (5.6, 5.7) -- (4.5, 5.7);
    \node[fill=boxmech,text=white,rounded corners=2pt,
          font=\small\bfseries] at (6.2, 5.7)
        {\SI{450}{\milli\metre} drop};
    \node[anchor=west,font=\footnotesize] at (7.2, 5.7)
        {$\sigma\approx\SI{523}{\kilo\pascal}$,\ $\Ke=\SI{221}{\joule}$};

    % Brace showing apparatus range
    \draw[decorate,decoration={brace,amplitude=6pt},gsvaccent,thick]
        (5.0,0.75) -- (5.0,5.7)
        node[midway,xshift=-0.4cm,rotate=90,font=\footnotesize\bfseries,
             text=gsvaccent] {Apparatus};
\end{tikzpicture}
\caption{Apparatus calibration to the Indian Railways field-loading
         envelope. The two drop heights span the in-service stress
         window from ordinary mainline traffic
         (\(\sim\!\SI{275}{\kilo\pascal}\), \SI{147}{\joule}) to
         heavy-haul / transition-zone conditions
         (\(\sim\!\SI{523}{\kilo\pascal}\), \SI{221}{\joule}).}
\label{fig:meth-calibration-ladder}
\end{figure}

The adjustability of the drop-height mechanism over the full
\SIrange{200}{450}{\milli\metre} range leaves open the possibility of finer
calibration to specific axle-load classes in future work.

\section{Sample preparation and test protocol}
\label{sec:meth-protocol}

\subsection{Specimen assembly}
\label{sec:meth-assembly}

The nine test specimens were assembled using a single standardised
protocol that produces reproducible specimens across the three pad
conditions and the two base conditions. The protocol comprises eleven
discrete steps, summarised here in the order of execution:
\begin{enumerate}[leftmargin=*,itemsep=2pt]
    \item Inspection and cleaning of the cylindrical test mould.
    \item Base preparation.
    \item Pad placement, where applicable, with fresh discs cut for every
          test.
    \item Tracer-particle placement at pre-assigned positions in the
          lower-middle of the ballast fill, with no tracer-to-tracer or
          tracer-to-wall direct contacts.
    \item Bulk-fill placement of the remainder of the \SI{20}{\kilogram}
          specimen mass in two lifts.
    \item Final levelling of the top surface with the load-transfer
          bearing-plate verification.
    \item Load-cell zeroing and digital-indicator verification.
    \item Photographic baseline record.
    \item Hammer positioning and locking at the specified drop height.
    \item Blow delivery.
    \item Test termination, tracer recovery, and bulk-fill collection for
          sieve analysis.
\end{enumerate}

\begin{figure}[ht]
    \centering
    \includegraphics[width=0.55\textwidth]{image19.png}
    \caption{Assembling the test specimen in the mould (photograph without
             load-cell assembly).}
    \label{fig:meth-assembly}
\end{figure}

\subsection{Blow-by-blow test protocol}
\label{sec:meth-blow-protocol}

Each of the nine tests was conducted using an identical blow-by-blow
protocol. Twenty (20) free-fall blows were delivered per test from the
nominated drop height (either \SI{300}{\milli\metre} or
\SI{450}{\milli\metre}). The choice of 20 blows is guided by two
considerations. First, it is consistent with the number of blows used in
comparable published drop-weight studies on ballast
\citep{Nimbalkar2012, Koohmishi2020, Zhang2023}. Second, preliminary trials
showed that the settlement--blow curve flattens substantially between
blow~15 and blow~20 and enters a near-stable zone beyond blow~15; twenty
blows therefore captures both the transient settlement regime and the
early approach to the steady state.

The twenty blows were delivered in rapid succession, with an interval of
\SIrange{5}{8}{\second} between consecutive blows---long enough for the
operators to safely re-raise and re-lock the hammer, and to record the
peak-force reading, but short enough that cumulative heating of the rubber
or polyurethane pad is negligible. The hammer release was triggered by
two-operator manual action at the upper crosshead. After each blow the
hammer was secured at the crosshead before the peak-force value was read.
Axial settlement was measured at five blow milestones ($N=0, 5, 10, 15,
20$) with a depth gauge seated on the mould rim. The digital peak-force
indicator was observed and recorded for each test.

\subsection{Test matrix and nomenclature}
\label{sec:meth-matrix}

The nine test configurations are defined by the factorial crossing of the
three pad conditions (\gls{UR} / \gls{RUBM} / \gls{PUBM}), the two base
conditions (\gls{FB} / \gls{RB}), and the two impact energies
(\SI{300}{\milli\metre} and \SI{450}{\milli\metre} drop heights of the
\SI{50}{\kilogram} hammer, corresponding to nominal contact stresses of
\SI{275}{\kilo\pascal} and \SI{523}{\kilo\pascal} as derived in
\cref{sec:meth-calibration}). The nine cells listed in
\cref{tab:meth-matrix} cover the factorial with the three
\gls{FB}--\SI{450}{\milli\metre} combinations omitted because the apparatus
commissioning revealed that at the \SI{450}{\milli\metre} drop height the
sand cushion under the \gls{FB} condition would progressively lose bearing
capacity through cumulative compaction over twenty blows and would thereby
confound the pad-effect signal with a base-consolidation signal.

\begin{figure}[ht]
\centering
\begin{tikzpicture}[node distance=0.15cm,>=Stealth]
    % Column header row
    \node[labelrow] (hdr-fb300) at (0, 2.4)
        {FB -- \SI{300}{\milli\metre}\\\footnotesize\SI{275}{\kilo\pascal}, \SI{147}{\joule}};
    \node[labelrow] (hdr-rb300) at (3.4, 2.4)
        {RB -- \SI{300}{\milli\metre}\\\footnotesize\SI{275}{\kilo\pascal}, \SI{147}{\joule}};
    \node[labelrow] (hdr-rb450) at (6.8, 2.4)
        {RB -- \SI{450}{\milli\metre}\\\footnotesize\SI{523}{\kilo\pascal}, \SI{221}{\joule}};

    % Row labels
    \node[labelrow,anchor=east] (lbl-pubm) at (-1.7, 1.2) {PUBM};
    \node[labelrow,anchor=east] (lbl-rubm) at (-1.7, 0.0) {RUBM};
    \node[labelrow,anchor=east] (lbl-ur)   at (-1.7,-1.2) {UR};

    % Row 1: PUBM (green)
    \node[cellpubm] at (0, 1.2)   {\testlabel{T1}\\PUBM--FB--300};
    \node[cellpubm] at (3.4, 1.2) {\testlabel{T9}\\PUBM--RB--300};
    \node[cellpubm] at (6.8, 1.2) {\testlabel{T6}\\PUBM--RB--450};
    % Row 2: RUBM (orange)
    \node[cellrubm] at (0, 0.0)   {\testlabel{T2}\\RUBM--FB--300};
    \node[cellrubm] at (3.4, 0.0) {\testlabel{T8}\\RUBM--RB--300};
    \node[cellrubm] at (6.8, 0.0) {\testlabel{T5}\\RUBM--RB--450};
    % Row 3: UR (red)
    \node[cellur] at (0,  -1.2) {\testlabel{T3}\\UR--FB--300};
    \node[cellur] at (3.4,-1.2) {\testlabel{T4}\\UR--RB--300};
    \node[cellur] at (6.8,-1.2) {\testlabel{T7}\\UR--RB--450};
\end{tikzpicture}
\caption{The nine-cell factorial test matrix. Rows are pad conditions
         (\gls{PUBM} / \gls{RUBM} / \gls{UR}); columns are
         base--energy combinations. The three
         \gls{FB}--\SI{450}{\milli\metre} cells are excluded by design
         because the sand cushion would progressively lose bearing
         capacity at the higher drop energy (see text).}
\label{fig:meth-matrix}
\end{figure}

\begin{table}[ht]
\centering
\caption{The nine test configurations of the experimental programme.
         Systematic labels are \texttt{PAD--BASE--DROP}; the legacy
         laboratory labels (\acrshort{LH}, \acrshort{HH}) are listed for
         cross-reference with the raw laboratory records.}
\label{tab:meth-matrix}
\rowcolors{2}{rowbg}{white}
\begin{tabular}{@{}cclcccc@{}}
    \toprule
    \rowcolor{tablehdr}
    {\color{white}\textbf{Test}}
    & {\color{white}\textbf{Systematic label}}
    & {\color{white}\textbf{Legacy label}}
    & {\color{white}$\boldsymbol{h}$~(\si{\milli\metre})}
    & {\color{white}$\boldsymbol{\sigma}$~(\si{\kilo\pascal})}
    & {\color{white}$\boldsymbol{\Ke}$~(\si{\joule})}
    & {\color{white}\textbf{Peak~$F$~(\si{\kilo\newton})}}\\
    \midrule
    \testlabel{T1} & PUBM--FB--300 & PUBM-FB-LH & 300 & 275 & 147 & 19.4\\
    \testlabel{T2} & RUBM--FB--300 & RUBM-FB-LH & 300 & 275 & 147 & 19.4\\
    \testlabel{T3} & UR--FB--300   & UR-FB-LH   & 300 & 275 & 147 & 19.4\\
    \testlabel{T4} & UR--RB--300   & UR-RB-LH   & 300 & 275 & 147 & 19.4\\
    \testlabel{T5} & RUBM--RB--450 & RUBM-RB-HH & 450 & 523 & 221 & 37.0\\
    \testlabel{T6} & PUBM--RB--450 & PUBM-RB-HH & 450 & 523 & 221 & 37.0\\
    \testlabel{T7} & UR--RB--450   & UR-RB-HH   & 450 & 523 & 221 & 37.0\\
    \testlabel{T8} & RUBM--RB--300 & RUBM-RB-LH & 300 & 275 & 147 & 19.4\\
    \testlabel{T9} & PUBM--RB--300 & PUBM-RB-LH & 300 & 275 & 147 & 19.4\\
    \bottomrule
\end{tabular}
\end{table}

\subsection{Post-test sieve analysis}
\label{sec:meth-sieve}

At the end of each test the full bulk fill, together with any fines
recovered from the mould floor after tracer extraction, was sieved through
the same sieve series used for the pre-test gradation characterisation
(63, 50, 40, 31.5, 25, 20, 16, 12.5, 10, 4.75, 2.36, 0.6, 0.15 and
\SI{0.075}{\milli\metre}).

\begin{figure}[ht]
    \centering
    \includegraphics[width=0.55\textwidth]{image20.jpeg}
    \caption{Post-impact-test sieve analysis.}
    \label{fig:meth-sieve}
\end{figure}

\subsection{Bulk degradation indices}
\label{sec:meth-indices}

A single fouling-based degradation index is computed on the post-test
gradation---the Indian Railway Fouling Index \FIIN.

\paragraph{Indian Railway Fouling Index.}
Following the adaptation by \citet{Anbazhagan2012} of \citet{Selig1994}:
\begin{equation}
    \FIIN \;=\; \PoneM \;+\; \PtwoM,
    \label{eq:meth-FIIN}
\end{equation}
where \PoneM{} and \PtwoM{} are the weight percent passing the
\SI{0.075}{\milli\metre} and \SI{20}{\milli\metre} sieves respectively. The
first term captures fines; the second the total fraction lost from the
acceptable \gls{RDSO} ballast gradation.

\section{Image-based three-dimensional morphological analysis}
\label{sec:meth-morphology}

\subsection{Scan protocol and mesh quality control}
\label{sec:meth-scan-protocol}

The seven-step scan protocol applied identically to every particle in the
programme---both pre-test and post-test---comprises:
\begin{enumerate}[leftmargin=*,itemsep=2pt]
    \item Cleaning (water wash, soft-brush scrub, distilled-water rinse,
          oven-dry at \SI{105}{\DegC}).
    \item Mounting on a low-tack dark putty pad at the centre of the
          automatic turntable.
    \item Calibration and ambient-light verification.
    \item Multi-orientation scanning to ensure complete surface coverage.
    \item Mesh reconstruction.
    \item Four-check quality-control gate (watertightness, single body,
          mesh volume within the
          \SIrange{10000}{250000}{\milli\metre\cubed} range).
    \item Export of the validated mesh as a binary \gls{STL} file.
\end{enumerate}

\begin{figure}[ht]
    \centering
    \begin{subfigure}{0.45\textwidth}
        \centering
        \includegraphics[width=\linewidth]{image21.jpeg}
        \caption{Particle scanning process.}
    \end{subfigure}\hfill
    \begin{subfigure}{0.45\textwidth}
        \centering
        \includegraphics[width=\linewidth]{image22.jpeg}
        \caption{Ballast particle in EXScan software.}
    \end{subfigure}
    \caption{Three-dimensional particle scanning workflow used in the
             present study.}
    \label{fig:meth-scan}
\end{figure}

\begin{figure}[ht]
    \centering
    \includegraphics[width=0.55\textwidth]{image23.jpeg}
    \caption{Particle aligning for multiple scans to ensure complete
             surface coverage.}
    \label{fig:meth-aligning}
\end{figure}

\subsection{Twenty-two morphological descriptors}
\label{sec:meth-descriptors}

Twenty-two morphological descriptors are computed per mesh in three
families: nine size descriptors ($V$, $A$, $\Deq$, $S$, $I$, $L$, $\Vch$,
$\Ach$, and the three ellipsoid semi-axes $a$, $b$, $c$); five form
descriptors (flatness $F=S/I$, elongation $E=I/L$, aspect ratio $L/S$,
flatness index $\mathit{FI}=(L+I)/(2S)$, and Zingg-triangle position); and
eight roundness and texture descriptors (sphericity $\Psi$, convexity
$C=V/\Vch$, scanner-roundness $R$, angularity $\mathrm{Ang}$,
surface-texture index $T=A/\Ach$, sphericity ratio $\kappa_{s}=\Deq/L$, and
the cross-check pair $\nu_{\mathrm{ch}}$ and $A_{\text{ratio}}$). The
classical Wadell sphericity \citep{Wadell1932} and the Krumbein roundness
concept \citep{Krumbein1941} remain the two principal reporting axes for
grain-scale degradation.

\subsection{Python script}
\label{sec:meth-python}

The descriptors are computed by a purpose-built Python~3 pipeline based on
the \texttt{trimesh} \citep{DawsonHaggerty2019}, \texttt{NumPy}
\citep{Harris2020}, and \texttt{SciPy} \citep{Virtanen2020} libraries. The
pipeline executes seven sequential stages on each particle's
before-and-after mesh pair:
\begin{enumerate}[label=(\roman*),leftmargin=*]
    \item Load \gls{STL}.
    \item Run the four-check QC gate.
    \item Compute size descriptors.
    \item Compute form descriptors.
    \item Compute roundness and texture descriptors.
    \item Compute paired changes ($\Delta X$, $\%\Delta X$).
    \item Export to a multi-sheet Excel workbook with sheets for Before,
          After, Changes (After~$-$~Before), and Percent Changes.
\end{enumerate}

The complete script is reproduced in \cref{chap:appB} and has been
unit-tested against synthetic reference meshes (sphere, ellipsoid, cube).

\begin{figure}[ht]
    \centering
    \includegraphics[width=0.85\textwidth]{image24.jpeg}
    \caption{Screenshot of the running morphological-analysis script.}
    \label{fig:meth-script}
\end{figure}

\begin{chaptertakeaways}
\begin{itemize}[leftmargin=*,itemsep=2pt,topsep=2pt]
    \item The basalt ballast conforms to the \gls{RDSO} gradation
          envelope (Figure~\ref{fig:meth-psd}); the two under-ballast
          mats (\gls{PUBM}, \gls{RUBM}) are identified, measured, and
          instrumented in \cref{sec:meth-PUBM,sec:meth-RUBM}.
    \item The modified drop-weight apparatus delivers
          \SI{275}{\kilo\pascal} / \SI{147}{\joule} at
          \SI{300}{\milli\metre} drop and \SI{523}{\kilo\pascal} /
          \SI{221}{\joule} at \SI{450}{\milli\metre} drop, bracketing
          the Indian Railways field-loading envelope
          (Figure~\ref{fig:meth-calibration-ladder}).
    \item The nine-cell factorial test matrix
          (Figure~\ref{fig:meth-matrix}) crosses three pad conditions
          with two base--energy combinations; the three
          \gls{FB}--\SI{450}{\milli\metre} cells are excluded by
          design.
    \item Forty-five tracer particles (five per configuration) are
          scanned before and after impact and reduced to a
          twenty-two-descriptor per-particle record by the
          \texttt{ballast\_analyzer.py} pipeline (\cref{chap:appB}).
    \item Both methodological strands---bulk \FIIN\ and per-particle
          paired morphology---can now be executed and analysed
          (Chapters~\ref{chap:impact-results} and
          \ref{chap:morph-results}).
\end{itemize}
\end{chaptertakeaways}
